\section{Rhie--Chow 補正と CCD 高次精度の整合性}
\label{app:rc_precision}
% ═══════════════════════════════════════════════════════════════════════════════

式~\eqref{eq:rc-face} の補正項は $O(h^2)$ の誤差を持つが，CCD の $O(h^6)$ 精度を律速しない理由を示す．

フェイス $f$（$x_P$ と $x_E$ の中点，$x_f = (x_P+x_E)/2$，$x_E - x_P = h$）での各勾配評価の
テイラー展開：

\textbf{(1) 直接フェイス差分：}
$x_P = x_f - h/2$，$x_E = x_f + h/2$ として
\[
  p_E - p_P = h\,p'(x_f) + \frac{h^3}{24}p'''(x_f) + \Ord{h^5}
\]
を $h$ で割ると
\begin{equation}
  (\nabla p)_f \equiv \frac{p_E - p_P}{h} = p'(x_f) + \frac{h^2}{24}p'''(x_f) + \Ord{h^4}
  \label{eq:rc_face_taylor}
\end{equation}

\textbf{(2) CCD 勾配の算術平均：}
CCD は各セル中心で $O(h^6)$ 精度の1階微分を返す：
$(D^{(1)}_\mathrm{CCD}\,p)_P = p'(x_P) + \Ord{h^6}$，
$(D^{(1)}_\mathrm{CCD}\,p)_E = p'(x_E) + \Ord{h^6}$．
これらの算術平均をテイラー展開：
\[
  p'(x_P) = p'(x_f) - \tfrac{h}{2}p''(x_f) + \tfrac{h^2}{8}p'''(x_f) + \Ord{h^3},
  \quad
  p'(x_E) = p'(x_f) + \tfrac{h}{2}p''(x_f) + \tfrac{h^2}{8}p'''(x_f) + \Ord{h^3}
\]
\begin{equation}
  \overline{(\nabla p)}_f \equiv \frac{1}{2}\bigl[(D^{(1)}_\mathrm{CCD}\,p)_P + (D^{(1)}_\mathrm{CCD}\,p)_E\bigr]
    = p'(x_f) + \frac{h^2}{8}p'''(x_f) + \Ord{h^4}
  \label{eq:rc_avg_taylor}
\end{equation}
（奇数次項が相殺し，$O(h^2)$ の次数で $p'(x_f)$ を近似することに注意；$O(h^6)$ ではない）

\textbf{(3) Rhie--Chow 検出項（差）：}
\begin{equation}
  (\nabla p)_f - \overline{(\nabla p)}_f
  = \left(\frac{1}{24} - \frac{1}{8}\right)h^2 p'''(x_f) + \Ord{h^4}
  = -\frac{h^2}{12}p'''(x_f) + \Ord{h^4} = \Ord{h^2}
  \label{eq:rc_detection_taylor}
\end{equation}
補正量は $\Ord{\Delta t \cdot h^2}$．

\begin{itemize}
  \item \textbf{滑らかな圧力場：} 補正が $\Ord{h^2}$ であり，CSF モデル誤差 $\Ord{\varepsilon^2} \approx \Ord{h^2}$ と同程度．CCD の $\Ord{h^6}$ を律速する新たな誤差源とはならない．
  \item \textbf{チェッカーボード場（波長 $2h$）：} $p_E - p_P = \Ord{1}$，$\overline{(\nabla p)}_f \approx 0$ のため補正量が $\Ord{1}$ の大きな値となり非物理振動を除去する．
  \item \textbf{スペクトル選択的フィルタ：} 低波数（物理情報）は CCD $\Ord{h^6}$ で処理し，高波数（チェッカーボード）のみ $\Ord{h^2}$ で減衰する選択的フィルタとして機能する．
\end{itemize}

\textbf{(4) 表面張力 RC 補正の誤差（式~\eqref{eq:rc-face-balanced} の $(\bm{f}_\sigma)_f - \overline{(\bm{f}_\sigma)}_f$）：}
拡張 Rhie--Chow 式~\eqref{eq:rc-face-balanced} では，圧力補正項に加えて表面張力検出項
$(\bm{f}_\sigma^{n+1})_f - \overline{(\bm{f}_\sigma^{n+1})}_f$ が加わる．
フェイス直接差分 $(\bm{f}_\sigma)_f = (\kappa_f/We)(\psi_E-\psi_P)/h$ はテイラー展開から：
\[
  (\bm{f}_\sigma)_f = (\kappa/We)\nabla\psi(x_f) + \Ord{h^2}
\]
セル中心 CCD 平均 $\overline{(\bm{f}_\sigma)}_f$ も式~\eqref{eq:rc_avg_taylor} と同じ構造で：
\[
  \overline{(\bm{f}_\sigma)}_f = (\kappa/We)\nabla\psi(x_f) + \Ord{h^2}
\]
差は $(\bm{f}_\sigma)_f - \overline{(\bm{f}_\sigma)}_f = \Ord{h^2}$，補正量は $\Ord{\Delta t \cdot h^2}$．
したがって表面張力補正項も CCD の $\Ord{h^6}$ 精度を律速しない．

Balanced-Force 条件（§\ref{sec:balanced_force}）への影響はない：Rhie--Chow 補正はフェイス速度の安定化補助項であり，主方程式の CCD 演算子とは独立の役割を担う（詳細は §\ref{sec:rc_balanced_force} 参照）．

% ═══════════════════════════════════════════════════════════════════════════════
